\chapter{Developer guide}
\label{chap:dev-guide}
% Many code paths in this chapter: allow looser lines instead of overfull ones
\sloppy

This chapter is a practical guide for whoever continues the work on FlowSync: how to run each
part, how the code is organised, how a request travels through the system, which conventions to
follow, and how to extend the system in the most common ways. It complements the README files of
the repository, which keep the full list of commands.

\status{The guide reflects the code on the date of this box. Where a step depends on a part that
does not exist yet (the trained models, the frontend on the real backend),
the guide says so and describes the intended way of extending it, so that it can be followed when
the part is built.}

\lstset{basicstyle=\ttfamily\footnotesize, columns=fullflexible, keepspaces=true, frame=single,
        breaklines=true, xleftmargin=2pt, xrightmargin=2pt}

\section{Map of the repository}
\label{sec:dev-map}

\begin{table}[htpb]
  \centering
  \small
  \begin{tabular}{@{}p{4.4cm}p{9.4cm}@{}}
    \toprule
    \textbf{Path} & \textbf{Content} \\
    \midrule
    \texttt{README.md} & What FlowSync is, state of each part, how to start it, test accounts \\
    \texttt{CONTRATTI\_API.md} & The contracts A--D between the parts (reference for every API) \\
    \texttt{FRONTEND\_STRUTTURA.md} & Structure and choices of the frontend \\
    \texttt{src/frontend/} & React application (Vite, TypeScript) \\
    \texttt{src/backend/} & Spring Boot application, MySQL schema (Flyway), \texttt{compose.yaml} \\
    \texttt{src/ai/} & Python package \texttt{p06\_ai}: schema, generator, dispatcher, rules,
      teacher, dataset; \texttt{docs/} with design decisions and experiment log \\
    \texttt{Tesi/} & This document (LaTeX) \\
    \texttt{FE\_Mockups/}, \texttt{GANTT/}, \texttt{*.vpp} & Mockups, project plan, UML models \\
    \bottomrule
  \end{tabular}
  \caption{Map of the repository}
  \label{tab:repo-map}
\end{table}

\section{Running the parts}
\label{sec:dev-run}

Each part runs on its own; the frontend does not need the others thanks to its mock.

\paragraph{Frontend} (Node.js~22). It starts on \texttt{http://localhost:5173} with the mock of
the backend enabled by \texttt{VITE\_USE\_MOCKS=true} in \texttt{.env.development}.
\begin{lstlisting}
cd src/frontend
npm install
npm run dev            # development server with the mock
npm run build          # type check + production build
npm run lint           # oxlint;  npm run format = Prettier
\end{lstlisting}

\paragraph{Backend} (Java~21, Maven, Docker running). The profile \texttt{dev} starts MySQL from
\texttt{compose.yaml}, applies the migrations and loads the demonstration data. Swagger UI is at
\texttt{/swagger-ui.html}.
\begin{lstlisting}
cd src/backend
cp .env.example .env   # database password and JWT secret
mvn spring-boot:run -Pdev
mvn test               # integration tests (Testcontainers MySQL)
\end{lstlisting}

\paragraph{AI} (Python~3.11 or later). The commands of the data pipeline are modules of the
package; the data they produce go into \texttt{src/ai/data/}, which is not versioned.
\begin{lstlisting}
cd src/ai
python -m venv .venv && .venv/bin/python -m pip install -e ".[dev]"
PY=.venv/bin/python
$PY -m pytest -q
$PY -m p06_ai.generator --profile FACTORY --count 200 --seed 0 \
    --out data/generated
$PY -m p06_ai.baselines data/generated --csv data/baselines.csv
$PY -m p06_ai.teacher data/generated --out data/teacher --time-limit 10
$PY -m p06_ai.dataset --instances data/generated \
    --teacher data/teacher --out data/datasets --version v1
\end{lstlisting}
The AI service of contract~B, needed by the planning endpoints of the backend, runs with:
\begin{lstlisting}
.venv/bin/python -m pip install -e ".[service]"
.venv/bin/uvicorn p06_ai.service:app --port 8000
\end{lstlisting}
The backend finds it at \texttt{AI\_BASE\_URL} (default \texttt{http://localhost:8000}); if the
service is off, planning requests end as failed and the rest of the backend works.

\paragraph{Frontend on the real backend.} Set \texttt{VITE\_USE\_MOCKS=false} and start the
backend: the development server forwards \texttt{/api} to \texttt{localhost:8080}. A partial mode
is also possible: with the mock enabled, a call for which no mock handler exists goes through to
the real backend (with a warning in the console). Since the real backend accepts only its own
tokens, in partial mode the login must be real too, i.e.\ the mock handlers of authentication must
be removed first.

\section{How a request travels through the system}
\label{sec:dev-flow}

The planning request is the central flow and touches every part. Its steps, and whether each
step exists today, are the following.
\begin{enumerate}
  \item The page \emph{Nuova pianificazione} sends the text to the interpretation endpoint
        (\texttt{POST /plan-requests/interpret}) and shows the result (frontend: done;
        backend: done with rules in \texttt{InterpretService}; language model: to do).
  \item The page sends \texttt{POST /plan-requests}; \texttt{PlanningService} checks the request
        (422 with paths), saves it with its ProblemSpec and answers \texttt{202} with the id
        (done).
  \item After the commit, \texttt{PlanningRunner} (\texttt{@Async}) rebuilds the problem with
        \texttt{PlanningContextFactory} in working minutes (\texttt{WorkTimeScale}) and the
        \texttt{AiClient} calls \texttt{POST /optimize} three times, one per variant (done).
  \item The AI service validates the specification, completes it with M1, builds the plans with
        M2 inside the dispatcher, falls back to the rules when needed, and returns the PlanResult
        (service: done, \texttt{p06\_ai.service}; until the models exist it returns the best rule
        as a fallback).
  \item The backend calls the service twice more for the ATC and EDD baselines, computes the
        indicators (\texttt{PlanKpis}) and stores plans and assignments; an unreachable service
        ends the request as failed with an audit record (done).
  \item The frontend polls \texttt{GET /plan-requests/\{id\}} every 1.5~s until the request is no
        longer pending, then shows the three plans; the Gantt chart reads
        \texttt{GET /plans/\{id\}/gantt}, and a decision goes to
        \texttt{POST /plans/\{id\}/decisions} (frontend: done on the mock; backend: done).
\end{enumerate}
The other flows (decisions, executions, reports, rescheduling) follow the same pattern: the
frontend calls an endpoint of contract A, the backend applies the business rules and, for a new
plan, calls the AI service through contract B.

\section{Frontend internals}
\label{sec:dev-frontend}

\begin{itemize}
  \item \texttt{src/api/} is the only folder that knows URLs and formats: \texttt{client.ts}
        (token, errors in RFC~7807 format turned into \texttt{ApiError}), one file of functions
        per area, \texttt{types/} with the types of contract A, and \texttt{queryKeys.ts} with the
        keys of the cache.
  \item \texttt{src/features/<area>/} holds the pages of an area, its components and a
        \texttt{hooks.ts} that wraps the API functions in TanStack Query (\texttt{useQuery} to
        read, \texttt{useMutation} to write and then refresh the keys that changed).
  \item \texttt{src/mocks/} is the simulated backend: \texttt{data/} with the demonstration data
        and \texttt{handlers/} with one file per area, registered in \texttt{handlers/index.ts}.
        In the browser console, \texttt{sessionStorage.setItem('mock:fail','1')} makes every call
        fail with 503, to test the error states.
  \item \texttt{src/app/router.tsx} lists all the routes; \texttt{RequireAuth} protects a route
        with a permission declared in \texttt{features/auth/permissions.ts}.
  \item \texttt{src/components/} and \texttt{src/lib/} hold the shared pieces: page header, card,
        status tag, form field with label and error, formats of durations, dates and amounts,
        labels and colours of the states, and the rules of the Gantt chart
        (\texttt{lib/schedule.ts}).
\end{itemize}

\section{Backend internals}
\label{sec:dev-backend}

\begin{itemize}
  \item One package per module under \texttt{it.unicam.flowsync}
        (\texttt{security}, \texttt{configuration}, \texttt{orders}, \texttt{planning},
        \texttt{execution}, \texttt{audit}), each with the sub-packages \texttt{domain} (JPA
        entities), \texttt{repository}, \texttt{service} (business rules, transactions) and
        \texttt{web} (controllers and DTOs as Java records).
  \item All the paths start with \texttt{ApiPaths.V1} (\texttt{/api/v1}); permissions are
        declared on the methods, e.g.\ \texttt{@PreAuthorize("hasRole('ADMIN')")}.
  \item Errors are thrown as \texttt{ApiException} (\texttt{notFound}, \texttt{conflict},
        \texttt{forbidden}, \texttt{unprocessable} with a list of errors, each with a code and a
        path) and are turned into RFC~7807 responses in one place, so the frontend can show each
        message next to its field.
  \item The schema is created only by Flyway migrations in \texttt{db/migration}
        (\texttt{V1}, \texttt{V2}, \dots); Hibernate only validates it. An applied migration is
        never modified. Demonstration data live in \texttt{db/seed} (\texttt{V100} and above),
        loaded only by the \texttt{dev} profile.
  \item The planning module works in \emph{minutes of working time}: convert an instant with
        \texttt{WorkTimeScale.toMinute} (round up for starts and releases, down for due dates)
        and back with \texttt{toInstant} (the boundary between two shifts is the end of an
        activity, not a start). Activity ids outside the database are
        \texttt{<order>/<code>}, e.g.\ \texttt{C-101/T}.
  \item The rules of a manual move are in the pure class \texttt{PlanChangeValidator}; a new rule
        goes there, with a unit test, and also in the Gantt module of the frontend.
\end{itemize}

\section{AI internals}
\label{sec:dev-ai}

\begin{itemize}
  \item \texttt{schema.py} (Pydantic models of the instance) and \texttt{validation.py}
        (consistency checks with error codes);
  \item \texttt{converters/} (benchmarks into the schema), \texttt{generator.py} (profiles and
        generation), \texttt{calibration.py};
  \item \texttt{dispatcher.py}: the simulator, with \texttt{candidates()} and
        \texttt{apply(candidate)}; a \emph{rule} is any function
        \texttt{(Dispatcher, list[Candidate]) -> Candidate}, and \texttt{run(instance, rule,
        solver)} builds a whole plan with it;
  \item \texttt{rules.py} (EDD, ATC), \texttt{checker.py} (independent feasibility check),
        \texttt{plan.py} (PlanResult and indicators), \texttt{objective.py} (weighted cost);
  \item \texttt{teacher.py} (CP-SAT), \texttt{features.py} (features of the candidates) and
        \texttt{dataset.py} (labelled steps and split).
\end{itemize}

\section{Conventions}
\label{sec:dev-conventions}

\begin{itemize}
  \item \textbf{Contracts first.} A new field or endpoint is first written in
        \texttt{CONTRATTI\_API.md}, agreed with the other parts, and then implemented. The
        frontend may simulate it in the mock in the meantime, with exactly the agreed format.
  \item \textbf{One branch.} Work is committed and pushed directly to \texttt{main}; commit
        messages start with the area (\texttt{frontend:}, \texttt{AI:}, \texttt{docs:},
        \texttt{thesis:}, \dots).
  \item \textbf{Languages.} The user interface is in Italian; the public API uses
        \texttt{camelCase}, the AI \texttt{snake\_case}; this document is in English.
  \item \textbf{Checks before a commit.} Frontend: \texttt{npm run build}, \texttt{npm run lint},
        \texttt{npm run format:check}; backend: \texttt{mvn test}; AI: \texttt{pytest}.
  \item \textbf{Documentation.} A change that alters a contract, a command or the state of a part
        updates the corresponding README and the status boxes of this document.
\end{itemize}

\section{How to extend FlowSync}
\label{sec:dev-extend}

The recipes below cover the most common extensions. Each lists the files to touch, in order.

\paragraph{Add a field to an existing object} (e.g.\ a ``material'' on an order).
\begin{enumerate}
  \item Contract: add the field to the object in \texttt{CONTRATTI\_API.md}.
  \item Backend: a new migration \texttt{V<n>\_\_order\_material.sql} with the column; the field
        in the entity (\texttt{orders/domain}), in the request and response records
        (\texttt{orders/web/OrderDtos.java}) and in the service; a check in the integration test.
  \item Frontend: the field in \texttt{api/types/orders.ts}, in the mock data, and in the form
        (\texttt{features/orders/components/OrderForm.tsx}: Zod schema, default values, a
        \texttt{FormField}).
\end{enumerate}

\paragraph{Add a page with its endpoint.}
\begin{enumerate}
  \item Contract: the endpoint, its roles and its JSON.
  \item Frontend: type in \texttt{api/types/}, function in \texttt{api/<area>.ts}, key in
        \texttt{queryKeys.ts}, hook in \texttt{features/<area>/hooks.ts}, the page in
        \texttt{features/<area>/}, the route in \texttt{app/router.tsx} (inside
        \texttt{RequireAuth} with the right permission), the menu entry in
        \texttt{app/layouts/Sidebar.tsx}, and a mock handler so the page works before the
        backend.
  \item Backend: controller, service and repository in the package of the module, with
        \texttt{@PreAuthorize}, errors through \texttt{ApiException}, and an integration test.
  \item When the backend endpoint is ready, delete the mock handler: the page does not change.
\end{enumerate}

\paragraph{Connect an area of the frontend to the real backend.} Compare the real responses with
the types of the frontend (Swagger UI helps); fix the differences in the contract and on the side
that is wrong; then remove the handlers of that area from \texttt{mocks/handlers/index.ts} and
test the pages with \texttt{VITE\_USE\_MOCKS=false}. Small differences can also be absorbed in
\texttt{src/api/}, as done for the resources (\texttt{normalizeResource}).

\paragraph{Add a role or change a permission.} Backend: the role in the enumeration of the user
and in the \texttt{CHECK} constraint (a new migration), and the \texttt{@PreAuthorize} of the
methods. Frontend: the role in \texttt{api/types/auth.ts}, its label and its permissions in
\texttt{features/auth/permissions.ts}; menu and routes follow automatically.

\paragraph{Add a domain profile} (e.g.\ a hospital).
\begin{enumerate}
  \item Backend: a row in \texttt{domain\_profile} with code, name, default weights and active
        constraints (a migration, or a seed for demonstration data). The frontend shows it
        without changes.
  \item AI: a new \texttt{Profile} in \texttt{generator.py} (sizes, durations, skills, branches,
        due dates, setups), added to \texttt{PROFILES}; then generate instances, run the
        baselines and the teacher, and build a new dataset version, so that the models learn the
        new kind of problem.
\end{enumerate}

\paragraph{Add a dispatching rule.} Write a function \texttt{(Dispatcher, list[Candidate]) ->
Candidate} in \texttt{rules.py} (for instance SPT: the candidate with the smallest
\texttt{duration\_min}, with the same tie-breaking as EDD), a \texttt{solve\_<rule>} function that
calls \texttt{run}, a unit test that checks feasibility with \texttt{check\_plan}, and the rule in
the dictionary of \texttt{baselines.py} so that it is evaluated with the others.

\paragraph{Add a feature for the policy M2.} Add its name and description to \texttt{FEATURES} in
\texttt{features.py} and compute it in \texttt{candidate\_features}, keeping it independent of the
scale of the instance; then build a new dataset version. The same function serves training and
inference, so the model will see the feature in both.

\paragraph{Add a new kind of constraint.} This is the deepest change, because feasibility is
guaranteed in several places that must stay consistent: the schema and its validator, the
dispatcher (legal moves), the independent checker, the CP-SAT model of the teacher, and the
ProblemSpec built by the backend. Each needs a test with an instance where the constraint is
binding.

\paragraph{Publish a model version} (when the training exists). The pipeline will evaluate the
new version against ATC and EDD, store its artefact and register it in the table
\texttt{model\_version} (contract C) as a candidate; promoting it to active is a deliberate step,
after which the AI service loads it and the page \emph{Modello AI} shows it.
\fussy
